\section*{Chapitre \ref{ch:fm:case-studies-i} --- Études de cas I}

\begin{solution}{exo:fm:case-studies-i:1}
$2.3/(1-0.44)=4.11$ milliards de dollars à la fin de juillet. Le 28 septembre, les 3.6~milliards de dollars représentaient 90~\% de la valeur liquidative, si bien que les anciens investisseurs détenaient
$3.6/0.9-3.6=\$0.40$ milliard~: un rendement de $0.40/2.3-1=-82.6~\%$ depuis la fin d'août.
\end{solution}

\begin{solution}{exo:fm:case-studies-i:2}
Voir la \cref{def:fm:case-studies-i:liquidity}. En 1998, financement~: les prêteurs et les contreparties pouvaient appeler des marges plus vite que les fourchettes ne se résorbaient. En
2006, marché~: le fonds détenait la majeure partie des positions ouvertes sur certains contrats et n'avait personne à qui vendre. En août 2007, la liquidité de marché, via l'encombrement (crowding), s'est transformée
en tension de financement, par le jeu de l'effet de levier et des limites de risque.
\end{solution}

\begin{solution}{exo:fm:case-studies-i:3}
$100\,000/(0.1\times25\,000)=40$ jours.
\end{solution}

\begin{solution}{exo:fm:case-studies-i:4}
$1/28=3.6~\%$~; le bilan valait environ $28\times4.67=131$ milliards de dollars. Nets de couvertures, les positions de spread ont bougé bien moins que leur valeur brute un
jour normal, ce qui explique que l'effet de levier ait paru acceptable.
\end{solution}

\begin{solution}{exo:fm:case-studies-i:5}
Voir la preuve de la \cref{prop:fm:case-studies-i:days}. À 75~\% des positions ouvertes, le volume disponible pour la position est faible au regard de sa taille,
de sorte que tout taux de participation assez bas pour limiter l'impact représente des semaines d'exposition, et que les ventes du fonds constituent l'essentiel des ventes du marché.
\end{solution}

\begin{solution}{exo:fm:case-studies-i:6}
Jusqu'au cinquième jour, le détenteur conserve tout son portefeuille tandis que celui du vendeur se réduit~; il subit donc une plus grande part du mouvement de chaque jour. Une fois les ventes arrêtées,
l'impact temporaire se résorbe~: le détenteur récupère sur l'ensemble de son portefeuille, le vendeur seulement sur la moitié qu'il a conservée, puisqu'il a vendu le reste aux cours déplacés.
\end{solution}

\begin{solution}{exo:fm:case-studies-i:7}
Pour un recouvrement de 0.5 (0.52 mesuré), le vendeur perd 9.8~\% au vingtième jour et le détenteur 4.3~\%~; avec des portefeuilles sans lien, le vendeur perd 9.7~\% et le
détenteur gagne 0.1~\%. La perte du vendeur est son propre impact~; celle du détenteur provient du recouvrement.
\end{solution}

\begin{solution}{exo:fm:case-studies-i:8}
Une valeur en risque à un jour suppose que la position peut être sortie dans la journée à des cours proches de ceux du modèle. Une position qui demande des semaines pour être liquidée comporte des semaines
de risque de prix, et ses propres ventes font bouger le cours~; les positions du fonds gazier représentaient l'essentiel des positions ouvertes sur certains contrats.
\end{solution}

\begin{solution}{pb:fm:case-studies-i:1}
\begin{enumerate}
\item Le dossier public et les chiffres reconstitués à partir de celui-ci~; il évite les méchants, le recul rétrospectif et les mobiles que le dossier ne mentionne pas.
\item Voir la \cref{def:fm:case-studies-i:liquidity}.
\item Les pertes entraînent des appels de marge, les appels imposent des ventes, les ventes font bouger les cours, et les cours déplacés constituent des pertes pour tous ceux qui détiennent les mêmes positions
(\cref{fig:fm:case-studies-i:spiral}).
\item Voir la \cref{def:fm:case-studies-i:days}.
\item Voir l'\cref{dat:fm:case-studies-i:record}~: 4.67~milliards de dollars de \omterm{def:fm:the-proprietary-market-making-firm:ownfunds}{fonds propres} à la fin de 1997 pour un effet de levier de 28 pour 1, 2.3~milliards de dollars à la fin d'août après une
perte de 44~\%, et 3.6~milliards de dollars apportés par 14 sociétés le 28 septembre, soit 90~\% de la valeur liquidative.
\item Elle a jugé qu'une liquidation rapide des positions du fonds, ainsi que des positions liées d'autres acteurs, risquait de menacer des marchés déjà fragilisés.
\item Jusqu'à 100\,000 contrats en un mois, soit environ 40~\% des contrats NYMEX de l'hiver 2006--2007 et 75~\% de ceux de novembre~; plus de 2~milliards de dollars perdus
de la fin août à la mi-septembre, et le portefeuille de gaz naturel liquidé.
\item Des transactions qui se recouvrent, une volatilité et des corrélations plus élevées~; risque et effet de levier réduits~; 3~milliards de dollars de \omterm{def:fm:the-proprietary-market-making-firm:ownfunds}{fonds propres} apportés à un fonds quantitatif valant environ
3.6~milliards de dollars.
\item $-12.1~\%$ de janvier à juillet, $-44~\%$ en août, $-82.6~\%$ jusqu'au 28 septembre.
\item Les montants réels de 4.67, 4.11, 2.30 et 0.40 milliards de dollars~; à la moitié de l'effet de levier, 4.67, 4.39, 3.42 et 2.01 milliards de dollars.
\item Les mêmes rendements d'actifs, lorsqu'une partie de la baisse de septembre s'est produite en contrepartie d'un fonds en détresse~; aucun crédit pour des appels de marge moins nombreux~; et rien n'a été pris
sur les années précédentes, que la moitié de l'effet de levier aurait réduit de moitié.
\item 20 jours et 40 jours, pour un volume illustratif de 25\,000 contrats par jour.
\item Voir la \cref{prop:fm:case-studies-i:days}.
\item Vendeur $-5.4~\%$, $-13.6~\%$, $-9.7~\%$~; détenteur $-5.3~\%$, $-15.1~\%$, $-7.9~\%$~; le détenteur conserve davantage de son portefeuille pendant les ventes et donc davantage de la
reprise qui suit.
\item Un risque d'apparence modeste, pour une taille et un effet de levier qui faisaient de la sortie un événement de marché, et une dépendance envers les prêteurs, la contrepartie et d'autres fonds.
\item Associer l'effet de levier à la durée du financement, des limites sur la part des positions ouvertes et sur les \omterm{def:fm:case-studies-i:days}{jours de liquidation}, ainsi qu'un test de résistance à une sortie encombrée.
\item Estimer le recouvrement du portefeuille avec les positions publiques ou les mesures d'encombrement (livre~7), puis simuler le débouclage d'un grand vendeur qui se recouvre
avec le portefeuille, et lui imputer les mouvements qui en découlent.
\item L'effet de levier fixe jusqu'où les pertes atteignent les \omterm{def:fm:the-proprietary-market-making-firm:ownfunds}{fonds propres}~; la durée du financement fixe le temps dont dispose la société pour attendre leur retournement.
\item 2.01~milliards de dollars au lieu de 0.40~milliard à la moitié de l'effet de levier~; 20 jours à 20~\% du volume illustratif (40 à 10~\%)~; le vendeur perd 9.7~\% au vingtième jour
contre 7.9~\% pour le détenteur.
\item Associer chaque limite d'effet de levier à une durée de financement minimale, et plafonner pour chaque position les \omterm{def:fm:case-studies-i:days}{jours de liquidation} et la part des positions ouvertes. Soumettre le portefeuille à un test de résistance
à une sortie encombrée, et décider à l'avance s'il faut vendre ou conserver.
\end{enumerate}
\end{solution}

\begin{solution}{iq:fm:case-studies-i:1}
Financement~: un prêteur relève les décotes et la société ne peut pas honorer l'appel de marge. Marché~: une position si grande par rapport au volume que la vendre fait bouger le cours.
\iqlookfor{les deux définitions et la manière dont elles s'alimentent mutuellement.}
\end{solution}

\begin{solution}{iq:fm:case-studies-i:2}
$1/25=4~\%$. Les actifs bruts surestiment le risque des positions de spread couvertes~; la mesure pertinente est le risque net en scénario de stress, lorsque les spreads bougent ensemble.
\iqlookfor{brut contre net.}
\end{solution}

\begin{solution}{iq:fm:case-studies-i:3}
Par les \omterm{def:fm:case-studies-i:days}{jours de liquidation} et l'impact d'une sortie, et non par une mesure de risque à un jour~; limiter la part des positions ouvertes et les \omterm{def:fm:case-studies-i:days}{jours de liquidation}.
\iqlookfor{coût de sortie contre risque à un jour.}
\end{solution}

\begin{solution}{iq:fm:case-studies-i:4}
Cela dépend du financement et des limites~: si la perte est une sortie encombrée que la société peut tenir, conserver la position récupère davantage~; si un appel de marge ou une limite
forcera la réduction plus tard, réduire tôt coûte moins cher. Vérifier d'abord le recouvrement et l'horizon de financement.
\iqlookfor{la décision dépend du financement, et pas seulement du signal.}
\end{solution}

\begin{solution}{iq:fm:case-studies-i:5}
Une liquidation forcée ferait bouger les cours contre leurs propres positions dans les mêmes transactions, et contre leurs expositions au fonds.
\iqlookfor{les portefeuilles propres des créanciers.}
\end{solution}

\begin{solution}{iq:fm:case-studies-i:6}
Les positions du portefeuille, leurs volumes et volatilités, un modèle d'impact, ainsi que la taille et le recouvrement des portefeuilles des autres détenteurs, qu'il faut estimer à partir de
positions publiques, de mesures d'encombrement ou d'épisodes passés.
\iqlookfor{le recouvrement non observé.}
\end{solution}
